\section{半年単位で再分類する — Release TopicからExecution Responsibilityへ}
\label{analysis:half-year-reclassification}

月刊で追ってきた「モデル公開」「Agent」「マルチモーダル」「Runtime」「Safety」という話題は、
半年単位では別々の流行として並べるより、\textbf{誰が実行結果に責任を持つ層なのか}という軸で
再分類した方が、2025年後期の変化を説明しやすい。これは個々の一次資料に直接書かれた結論ではなく、
本号で検証済みの時系列を横断して行う編集上の再分類である。

\begin{center}
\small
\renewcommand{\arraystretch}{1.18}
\begin{tabularx}{\linewidth}{@{}p{0.20\linewidth}p{0.25\linewidth}X@{}}
\toprule
責任の層 & 半年で前景化した対象 & 再分類すると見える変化 \\
\midrule
Delivery & API、product、open weights &
model announcementだけでなく、API提供、product rollout、open-weight配布を別identityとして追う必要が増えた。
GPT-5とgpt-ossは、同じ「モデル」の話でもdelivery responsibilityが異なる。 \cite{src-05738f6c55b7,src-964121647555,src-12865e8612c7,src-1c46fdf17c79} \\
Execution & harness、context、tool、browser/computer use &
ChatGPT agent、Claude系agent platform、Gemini APIの更新は、能力を実行へ接続するharnessやtool surfaceを
model本体と分けて観察する必要を強めた。 \cite{src-b074ccdb2d9e,src-7a7213b68631,src-f6aa679babd7,src-a24fd97eb3d0} \\
Action space & video、world model、robotics &
Sora 2やDeepMindのworld-model / robotics研究は、出力生成だけでなく環境・行動空間との接続を
独立した観測軸にした。 \cite{src-661a3476d3d5,src-8530f5b0e216,src-dd5d9b8662f3} \\
Runtime & library、serving、integration &
vLLMやTransformersなどのrelease chronologyは、上流model releaseとdownstream supportが
同じ日付でも同じ意味でもないことを明確にした。 \cite{src-43eaa2f12960,src-58ef6e84aa65,src-455a78bc36fa,src-2595cdd4fabc} \\
Control & evaluation、safeguard、prompt-injection defense &
cross-lab safety evaluation、gpt-oss-safeguard、Atlas hardeningは、controlをrelease時の付帯説明ではなく
継続的なengineering responsibilityとして追う必要を示した。 \cite{src-ad13fe7d3321,src-2081c76ea2b4,src-0ad9fbef8c87} \\
\bottomrule
\end{tabularx}
\renewcommand{\arraystretch}{1.0}
\end{center}
\normalsize

この再分類で重要なのは、modelの重要性が下がったという意味ではない点である。
むしろmodel capabilityが増えるほど、delivery、execution、runtime、controlの各層が
「その能力をどこで、どの条件で、どの責任境界の下に使えるか」を決める比重が増した、と読む方が正確である。

\Needspace{0.38\textheight}
\section{Cross-month Comparison — 同じ責任を7月から12月まで追う}
\label{analysis:cross-month}

月ごとのheadlineを並べるだけでは、同じ責任が半年の中でどう変形したかが見えにくい。
そこで7--8月、9--10月、11--12月を同じ責任軸で比較する。

\begin{center}
\footnotesize
\renewcommand{\arraystretch}{1.18}
\begin{tabularx}{\linewidth}{@{}p{0.14\linewidth}p{0.27\linewidth}p{0.27\linewidth}X@{}}
\toprule
責任 & 7--8月 & 9--10月 & 11--12月 \\
\midrule
Execution &
ChatGPT agentがresearch、browser、code/tool executionを一つの実行面へ寄せ、
model外側のharnessを可視化した。 \cite{src-b074ccdb2d9e,src-7a7213b68631} &
Codex更新、AgentKit、browser/computer-use系の展開が重なり、tool接続と実行surfaceが増えた。
\cite{src-6deeebce1719,src-0661418088b8,src-a24fd97eb3d0} &
Claude 4.5系やcoding specializationの更新が続き、単一launchよりfamilyと実行環境の
継続更新を追う必要が強まった。 \cite{src-f6aa679babd7,src-55da04f42e32,src-cd8d6c8023f2} \\
Delivery / Model family &
gpt-ossとGPT-5がopen-weight deploymentとfrontier API/productという異なるdeliveryを同時期に示した。
\cite{src-12865e8612c7,src-1c46fdf17c79,src-05738f6c55b7} &
Qwen、DeepSeekなど複数familyの更新とservice availabilityが重なり、upstream releaseとprovider availabilityを
分ける必要が増えた。 \cite{src-eac8dfbeb78c,src-e547d0d8cbe3} &
GPT-5.1 / 5.2、Kimi、Mistral、GLMなどの更新が続き、観測単位が「一度の発表」から
family lifecycleへ移った。 \cite{src-5bea303cfdb1,src-10ab1c39be14,src-9191d43918fc,src-d8b7c1c36830,src-9156c37fb1c2} \\
Runtime &
vLLM v0.10、Transformers v4.54/4.55などが上流の変化をdownstreamへ取り込む時系列を作った。
\cite{src-43eaa2f12960,src-a783d282600b} &
vLLM v0.11やTransformers v4.57など、runtime/library側の更新がmodel churnと並行した。
\cite{src-58ef6e84aa65,src-2595cdd4fabc} &
vLLM v0.13、SGLang、TensorRT-LLMの更新が続き、serving側も半年を通じた独自cadenceを持った。
\cite{src-455a78bc36fa,src-ee0eb9562688,src-04b5c0f10b59} \\
Control &
GPT-5 safe-completionsとcross-lab safety evaluationが、releaseと評価を別のcontrol surfaceとして示した。
\cite{src-aec54e92a7f3,src-ad13fe7d3321} &
scheming研究、gpt-oss-safeguard、Atlas/OWLが、model safetyだけでなくagent/browser securityを
control対象へ広げた。 \cite{src-bef6cbe2a83e,src-2081c76ea2b4,src-240a05c02f85,src-3556c47e80db} &
Atlas hardeningは、公開後もprompt injection対策が継続することを示し、controlを
post-launch engineeringとして読む必要を強めた。 \cite{src-0ad9fbef8c87} \\
\bottomrule
\end{tabularx}
\renewcommand{\arraystretch}{1.0}
\end{center}
\normalsize

この比較から読めるのは「半年で一つの勝者へ収束した」ということではない。
むしろ同じ責任層が別々のcadenceで更新され、後半ほど層同士の依存関係が目立つようになった、という時系列である。

\Needspace{0.36\textheight}
\section{Cross-layer Synthesis — 層をまたぐ結合点}
\label{analysis:cross-layer}

個別章を横断すると、2025年後期の特徴は各layerの強化だけでなく、
\textbf{layer間の結合点がreader-facingな技術論点になったこと}にある。
以下は因果関係や性能優位を主張するものではなく、本号で確認した独立した時系列を
同じsystem boundary上へ配置した編集上のsynthesisである。

\paragraph{Agent $\times$ Runtime}
Agentの価値はmodel capabilityだけでは決まらず、tool call、context、browser/computer-useを実行へ接続するharnessと、
その下でmodel churnを吸収するruntime/libraryの両方に依存する。
ChatGPT agentやClaude系agent platformの時系列と、vLLM / Transformersの時系列を分離して追うことで、
「Agent公開」と「downstream runtime support」を同一eventとして扱わずに結合関係を読める。
\cite{src-b074ccdb2d9e,src-f6aa679babd7,src-43eaa2f12960,src-2595cdd4fabc}

\paragraph{Multimodal / Action space $\times$ Control}
video、world model、robotics、computer-useがaction spaceを広げるほど、
control boundaryはtext responseだけでは閉じない。Sora 2のreleaseとSystem Card、
DeepMindのworld-model / robotics research、Atlasのbrowser security hardeningは、
それぞれ別のidentityを保ちながら「何を生成できるか」と「どこへ作用できるか」を分けて読む必要を示す。
\cite{src-661a3476d3d5,src-8530f5b0e216,src-dd5d9b8662f3,src-0ad9fbef8c87}

\paragraph{Open Weight $\times$ Runtime / Delivery}
open weightsでは、upstreamのweight公開と、実際のlibrary / serving runtimeへのintegration、
利用者側deploymentが別責任になる。gpt-ossとTransformers / vLLMの時系列は、
「公開された」と「特定runtimeで扱える」を同義にしないための代表例である。
\cite{src-12865e8612c7,src-1c46fdf17c79,src-2595cdd4fabc,src-43eaa2f12960}

\begin{claimboundary}[横断分析の境界]
本節の「責任」「結合」「execution stack」という整理は編集上の推論であり、
引用資料が同一のsystem architectureや因果関係を共同で主張しているわけではない。
vendor / projectが報告する性能値や優位性も、独立再現された一般則へ変換しない。
ここでの目的は、model announcement、API/product availability、runtime integration、
safety/security controlを同一identityへ潰さず、半年の技術変化を比較できる観測軸へ整理することである。
\end{claimboundary}
